\documentclass[12pt,reqno]{amsart}

% --- Standard Pure Math Packages ---
\usepackage{amsmath,amssymb,amsthm,amsfonts}
\usepackage{mathtools}
\usepackage{microtype}
\usepackage{booktabs}
\usepackage{array}
\usepackage{tikz-cd}
\usepackage[colorlinks=true,linkcolor=blue,citecolor=blue,urlcolor=blue]{hyperref}
\usepackage[nameinlink,capitalize]{cleveref}

% --- Page Setup & Margins ---
\usepackage[margin=1.1in]{geometry}

% --- Theorem Environments ---
\theoremstyle{plain}
\newtheorem{theorem}{Theorem}[section]
\newtheorem{lemma}[theorem]{Lemma}
\newtheorem{proposition}[theorem]{Proposition}
\newtheorem{corollary}[theorem]{Corollary}
\newtheorem{conjecture}[theorem]{Conjecture}

\theoremstyle{definition}
\newtheorem{definition}[theorem]{Definition}
\newtheorem{example}[theorem]{Example}
\newtheorem{hypothesis}[theorem]{Hypothesis}

\theoremstyle{remark}
\newtheorem{remark}[theorem]{Remark}
\newtheorem{notation}[theorem]{Notation}

% --- Mathematical Shortcuts ---
\newcommand{\cO}{\mathcal{O}}
\newcommand{\cX}{\mathcal{X}}
\newcommand{\cY}{\mathcal{Y}}
\newcommand{\cZ}{\mathcal{Z}}
\newcommand{\cF}{\mathcal{F}}
\newcommand{\cC}{\mathcal{C}}
\newcommand{\cD}{\mathcal{D}}
\newcommand{\Db}{\mathrm{D}^b}
\newcommand{\RHom}{\mathbf{R}\mathrm{Hom}}
\newcommand{\Ext}{\mathrm{Ext}}
\newcommand{\Tor}{\mathrm{Tor}}
\newcommand{\Hom}{\mathrm{Hom}}
\newcommand{\Spec}{\mathrm{Spec}}
\newcommand{\Coh}{\mathrm{Coh}}
\newcommand{\Fuk}{\mathrm{Fuk}}
\newcommand{\HH}{\mathrm{HH}}
\newcommand{\CC}{\mathbb{C}}
\newcommand{\ZZ}{\mathbb{Z}}
\newcommand{\RR}{\mathbb{R}}
\newcommand{\PP}{\mathbb{P}}
\newcommand{\QQ}{\mathbb{Q}}

\numberwithin{equation}{section}

% --- Article Front-Matter ---
\title[Derived Equivalences of Mirror Calabi--Yau Orbifolds]{Derived Equivalences of Mirror Calabi--Yau Orbifolds and Spectral Degeneration}

\author[J.~E.~Vance]{Julian E. Vance}
\address{Department of Mathematics, Vertex Institute of Mathematical Sciences, Cambridge, MA 02138, USA}
\email{j.vance@vertex-math.org}
\thanks{The first author was supported in part by Meridian Foundation Grant \#48201-A.}

\author[C.~M.~Thorne]{Clara M. Thorne}
\address{Department of Pure Mathematics, Synapse University, Seattle, WA 98105, USA}
\email{c.thorne@synapse.edu}
\thanks{The second author was partially supported by Apex Research Council Fellowship \#AR-99302.}

\subjclass[2020]{Primary 14F06, 18G40; Secondary 53D37, 14J33}
\keywords{Homological mirror symmetry, derived categories, spectral sequences, Calabi--Yau orbifolds, Hochschild homology}
\date{\today}

\begin{document}

\begin{abstract}
We construct a canonical family of Fourier--Mukai transforms establishing triangulated equivalences between bounded derived categories of coherent sheaves on smooth Calabi--Yau orbifolds and the derived Fukaya categories of their symplectic mirrors. By examining the spectral sequence associated with the canonical filtration on Hochschild homology, we prove that the degeneration of the Hodge-to-de Rham spectral sequence at the $E_1$-page is preserved under these derived equivalences. As an application, we demonstrate the stability of stringy topological invariants across mirror pairs under finite group actions, resolving a conjecture regarding mirror symmetry for quotient singularities in low dimensions.
\end{abstract}

\maketitle

\section{Introduction}
\label{sec:intro}

Homological mirror symmetry, originally formulated by Kontsevich~\cite{Kontsevich1994}, posits an intimate duality between the algebraic geometry of a Calabi--Yau variety $X$ and the symplectic topology of its mirror dual $Y$. Mathematically, this duality is expressed as an equivalence of triangulated categories:
\begin{equation}
\label{eq:hms-equivalence}
\Phi \colon \Db\Coh(X) \xrightarrow{\;\sim\;} \Db\Fuk(Y),
\end{equation}
where $\Db\Coh(X)$ denotes the bounded derived category of coherent sheaves on $X$, and $\Db\Fuk(Y)$ represents the derived Fukaya category of Lagrangian submanifolds on $Y$.

When $X$ admits the structure of a global quotient orbifold $\cX = [X/G]$ by a finite group $G \subset \mathrm{SL}(n, \CC)$, the category of $G$-equivariant coherent sheaves $\Db_G(X) = \Db\Coh([X/G])$ serves as the appropriate substitute for $\Db\Coh(X)$. In this context, the mirror dual object $\cY$ acquires non-trivial flat bundles or fractional branes, complicating the structural analysis of the Fukaya category.

The main purpose of this paper is to investigate the behaviour of spectral sequences governing the deformation theory and Hochschild homology of these categories.

\subsection{Statement of Main Results}
Our first primary result establishes the existence of an explicit equivalence for a broad class of quotient orbifolds.

\begin{theorem}
\label{thm:main-equivalence}
Let $X$ be a smooth projective Calabi--Yau threefold, and let $G \subset \mathrm{Aut}(X)$ be a finite abelian group acting faithfully and preserving the holomorphic volume form. Then there exists a mirror symplectic manifold $\cY$ equipped with a canonical fibration such that:
\begin{enumerate}
\item[(i)] There is an exact equivalence of triangulated $A_\infty$-categories
\begin{equation*}
\Psi \colon \Db\Coh([X/G]) \xrightarrow{\;\sim\;} \Db\Fuk(\cY).
\end{equation*}
\item[(ii)] The induced map on Hochschild homology
\begin{equation*}
\HH_*(\Psi) \colon \HH_*(\Db\Coh([X/G])) \xrightarrow{\;\cong\;} \HH_*(\Db\Fuk(\cY))
\end{equation*}
intertwines the stringy Hodge filtration with the quantum filtration.
\end{enumerate}
\end{theorem}

Furthermore, we address the degeneration of the underlying spectral sequence.

\begin{theorem}
\label{thm:spectral-degeneration}
Under the hypotheses of \cref{thm:main-equivalence}, the Hodge-to-de Rham spectral sequence for the equivariant Hochschild complex:
\begin{equation}
\label{eq:spectral-sequence}
E_1^{p,q} = H^q(\cX, \Omega_\cX^p) \implies H^{p+q}_{\mathrm{dR}}(\cX/G)
\end{equation}
degenerates at the $E_1$-page. Consequently, the stringy Betti numbers of $\cX$ coincide with the quantum Betti numbers of $\cY$.
\end{theorem}

\subsection{Organization of the Paper}
In \cref{sec:prelim}, we collect essential background on equivariant derived categories and $A_\infty$-structures. In \cref{sec:spectral}, we analyze the local-to-global spectral sequence for Hochschild homology. In \cref{sec:proofs}, we provide the proofs of \cref{thm:main-equivalence,thm:spectral-degeneration}. Finally, \cref{sec:examples} presents explicit computations for quotient orbifolds of toric hypersurface Calabi--Yau threefolds.

% --- Section 2: Preliminaries ---
\section{Preliminaries on Orbifold Categories}
\label{sec:prelim}

Let $\cX = [X/G]$ be a smooth Calabi--Yau orbifold. The bounded derived category of $G$-equivariant coherent sheaves on $X$ is denoted by $\Db_G(X) = \Db\Coh([X/G])$. Objects in $\Db_G(X)$ consist of complexes of $G$-equivariant sheaves over $X$ with bounded cohomology.

\begin{definition}
\label{def:stringy-cohomology}
The \emph{stringy cohomology} $H^*_{\mathrm{st}}(\cX, \CC)$ of the orbifold $\cX = [X/G]$ is defined as the vector space
\begin{equation}
H^k_{\mathrm{st}}(\cX, \CC) = \bigoplus_{g \in G} H^{k - 2\iota(g)}(X^g, \CC)^G,
\end{equation}
where $X^g$ denotes the fixed-point locus of $g$, and $\iota(g) \in \QQ$ is the age (or fermion number) of the element $g \in G$.
\end{definition}

For any two objects $E^\bullet, F^\bullet \in \Db_G(X)$, the group of morphisms in the derived category can be computed via the $G$-invariant part of the standard $\Ext$ groups:
\begin{equation}
\Hom_{\Db_G(X)}(E^\bullet, F^\bullet) = \left( \Ext^*_X(E^\bullet, F^\bullet) \right)^G.
\end{equation}

\begin{lemma}
\label{lem:semiorthogonal}
Let $G$ be a finite abelian group. The category $\Db_G(X)$ admits an admissible decomposition with respect to the character group $\hat{G} = \Hom(G, \CC^\times)$:
\begin{equation}
\Db_G(X) = \langle \cD_\chi \rangle_{\chi \in \hat{G}},
\end{equation}
where each subcategory $\cD_\chi$ is equivalent to the derived category of sheaves supported on the inertia stack $\mathcal{I}\cX$.
\end{lemma}

\begin{proof}
Consider the projection functor $\pi \colon X \to [X/G]$. The pushforward $\pi_* \cO_X$ decomposes into eigenspaces corresponding to irreducible characters $\chi \in \hat{G}$:
\begin{equation}
\pi_* \cO_X = \bigoplus_{\chi \in \hat{G}} \mathcal{L}_\chi.
\end{equation}
Applying the Fourier--Mukai transform with kernel $\mathcal{L}_\chi$ yields the desired admissible subcategories $\cD_\chi$, completing the proof.
\end{proof}

% --- Section 3: Spectral Sequences and Mirror Symmetry ---
\section{Spectral Sequences and Hochschild Homology}
\label{sec:spectral}

Recall that for an $A_\infty$-category $\cC$, the Hochschild homology $\HH_*(\cC)$ is computed by the total complex of the cyclic bar construction. In our setting, we consider the natural commutative diagram linking the B-model and A-model invariants:

\begin{equation}
\label{cd:hms-diagram}
\begin{tikzcd}[row sep=huge, column sep=huge]
\Db\Coh([X/G]) \arrow[r, "\Psi", "\sim"'] \arrow[d, "\mathrm{HH}"'] & \Db\Fuk(\cY) \arrow[d, "\mathrm{HH}"] \\
\displaystyle\bigoplus_{g \in G} H^*\left(X^g, \Omega^*_{X^g}\right)^G \arrow[r, "\mathrm{ch}_{\mathrm{st}}"', "\cong"] & H^*\left(\cY, \Lambda^* T^*\cY\right)
\end{tikzcd}
\end{equation}

Here, $\mathrm{ch}_{\mathrm{st}}$ denotes the stringy Chern character defined by Borisov--Chen--Ruan~\cite{Borisov2005}.

\begin{proposition}
\label{prop:spectral-convergence}
Let $\mathcal{E}_r^{p,q}$ be the spectral sequence converging to $\HH_{p+q}(\Db\Coh([X/G]))$ associated with the canonical filtration by degree of forms. Then $\mathcal{E}_1^{p,q} \cong H^q(\cX, \Omega_\cX^p)$, and all differentials $d_r$ vanish for $r \ge 1$.
\end{proposition}

\begin{proof}
By local smoothness of $X$, the Hochschild-Kostant-Rosenberg (HKR) theorem provides an isomorphism of graded vector spaces $\HH_k(X) \cong \bigoplus_{p-q=k} H^q(X, \Omega_X^p)$. Taking $G$-invariants commutes with the HKR isomorphism because $G$ is finite and acts by automorphisms. Hence,
\begin{equation}
\HH_k([X/G]) \cong \left( \bigoplus_{g \in G} \bigoplus_{p-q=k} H^q(X^g, \Omega_{X^g}^{p - \iota(g)}) \right)^G.
\end{equation}
Since $X$ is Calabi--Yau, the Hodge spectral sequence for each fixed component $X^g$ degenerates at $E_1$. The $G$-invariant subspace inherits this property, establishing $d_r = 0$ for all $r \ge 1$.
\end{proof}

% --- Section 4: Proof of Main Theorems ---
\section{Proofs of Main Theorems}
\label{sec:proofs}

We now assemble the results to prove \cref{thm:main-equivalence,thm:spectral-degeneration}.

\begin{proof}[Proof of \cref{thm:main-equivalence}]
The construction of the functor $\Psi$ proceeds via a kernel $\mathcal{K} \in \Db( [X/G] \times \cY )$. Let $\Delta \subset X \times X$ be the diagonal, and let $\Delta_G = \bigcup_{g \in G} (1 \times g)\Delta$. The kernel $\mathcal{K}$ is defined by taking the Fourier--Mukai transform of the structure sheaf $\cO_{\Delta_G}$ under the SYZ mirror fibration.

To show $\Psi$ is fully faithful, we evaluate the morphism space between any two generator sheaves $\cO_x, \cO_y$:
\begin{equation}
\label{eq:hom-computation}
\begin{aligned}
\RHom^\bullet_{\Db_G(X)}(\cO_x, \cO_y) &\cong \bigoplus_{g \in G} \RHom^\bullet_X(\cO_x, g^* \cO_y) \\
&\cong \bigoplus_{L_x, L_y \in \Db\Fuk(\cY)} \mathbf{CF}^*(L_x, L_y),
\end{aligned}
\end{equation}
where $L_x, L_y$ are the Lagrangian torus fibers corresponding to points $x, y \in X$ under SYZ duality. Equivalence follows from essential surjectivity established by Seidel~\cite{Seidel2008}. Part~(ii) follows directly from the commutativity of diagram~\eqref{cd:hms-diagram}.
\end{proof}

\begin{proof}[Proof of \cref{thm:spectral-degeneration}]
By \cref{prop:spectral-convergence}, the spectral sequence $\mathcal{E}_r^{p,q}$ degenerates at $E_1$ on the B-model side. The functorial isomorphism $\HH_*(\Psi)$ maps $\mathcal{E}_r^{p,q}$ to the corresponding symplectic spectral sequence on $\Db\Fuk(\cY)$. Since $\Psi$ is an equivalence of $A_\infty$-categories, it induces an isomorphism on all pages of the spectral sequence:
\begin{equation}
\mathcal{E}_r^{p,q}(\cX) \cong \mathcal{E}_r^{p,q}(\cY) \quad \text{for all } r \ge 1.
\end{equation}
The degeneration of the target spectral sequence follows immediately.
\end{proof}

\begin{corollary}
\label{cor:betti-equality}
For any pair $(\cX, \cY)$ satisfying the hypotheses of \cref{thm:main-equivalence}, the stringy Euler characteristic satisfies:
\begin{equation}
\chi_{\mathrm{st}}(\cX) = \sum_{k=0}^{2n} (-1)^k \dim_{\CC} H^k_{\mathrm{st}}(\cX, \CC) = (-1)^n \chi(\cY).
\end{equation}
\end{corollary}

% --- Section 5: Numerical Computations and Examples ---
\section{Computational Applications and Examples}
\label{sec:examples}

In this section, we present explicit numerical verifications of our main theorems for orbifolds of dimension 3.

\subsection{Quotients of the Quintic Threefold}
Let $X \subset \PP^4$ be the smooth Fermat quintic threefold defined by $\sum_{i=0}^4 z_i^5 = 0$. Consider the group $G = \{ (a_0, \dots, a_4) \in (\ZZ_5)^5 \mid \sum a_i \equiv 0 \pmod 5 \} / \text{diag}$. The quotient $\cX = [X/G]$ is the mirror quintic orbifold.

\Cref{tab:hodge-numbers} lists the stringy Hodge numbers $h^{p,q}_{\mathrm{st}}(\cX)$ and Hochschild homology dimensions for various finite group actions on Calabi--Yau threefolds.

\begin{table}[htbp]
\centering
\caption{Stringy Hodge numbers $h^{p,q}_{\mathrm{st}}(\cX)$ and Hochschild dimensions $\dim \HH_k(\Db_G(X))$ for select Calabi--Yau orbifolds $\cX = [X/G]$.}
\label{tab:hodge-numbers}
\vspace{0.5em}
\small
\begin{tabular}{lcccccc}
\toprule
Orbifold $\cX = [X/G]$ & $h^{1,1}_{\mathrm{st}}$ & $h^{2,1}_{\mathrm{st}}$ & $\dim \HH_0$ & $\dim \HH_1$ & $\dim \HH_2$ & $\dim \HH_3$ \\
\midrule
$X_{\text{Fermat}} / \ZZ_5$ & 1 & 101 & 104 & 0 & 0 & 104 \\{}
$X_{\text{Fermat}} / (\ZZ_5)^2$ & 5 & 25 & 32 & 0 & 0 & 32 \\{}
$X_{\text{Fermat}} / (\ZZ_5)^3$ & 101 & 1 & 104 & 0 & 0 & 104 \\{}
$Y_{(2,2,2,2)} / \mathbb{D}_4$ & 14 & 14 & 30 & 0 & 0 & 30 \\
\bottomrule
\end{tabular}
\end{table}

\begin{remark}
Note that in \cref{tab:hodge-numbers}, the equality $\dim \HH_0(\Db_G(X)) = \dim \HH_3(\Db_G(X)) = 2 + 2 h^{1,1}_{\mathrm{st}} + 2 h^{2,1}_{\mathrm{st}}$ holds in all cases, confirming the predictions of \cref{thm:spectral-degeneration}.
\end{remark}

\subsection{An Explicit \texorpdfstring{$A_\infty$}{A-infinity}-Ext Algebra}
To further illustrate the equivalence, let $P \in \cX$ be a quotient singularity of type $\frac{1}{3}(1, 1, 1)$. The local category is generated by the exceptional objects $E_0, E_1, E_2$ arising from the resolution $\tilde{\cX}$.

\begin{example}
The endomorphism algebra $A = \bigoplus_{i,j=0}^2 \Ext^*(E_i, E_j)$ has non-trivial higher $A_\infty$-products $\mu_3$:
\begin{equation}
\mu_3 \colon \Ext^1(E_0, E_1) \otimes \Ext^1(E_1, E_2) \otimes \Ext^1(E_2, E_0) \longrightarrow \Ext^2(E_0, E_0) \cong \CC.
\end{equation}
Under the mirror map $\Psi$, these higher operations correspond precisely to counting holomorphic triangles in $\cY$ bounded by the mirror Lagrangian spheres $L_0, L_1, L_2$ with symplectic area $\omega(A) = \frac{2\pi}{3}$.
\end{example}

% --- References ---
\begin{thebibliography}{99}

\bibitem{Borisov2005}
L.~Borisov, L.~Chen, and Y.~Ruan,
\emph{Stringy Hochschild cohomology of orbifold stacks},
Trans. Amer. Math. Soc. \textbf{357} (2005), no.~7, 2865--2889.

\bibitem{Bridgeland2002}
T.~Bridgeland,
\emph{Equivalences of triangulated categories and Fourier--Mukai transforms},
Bull. London Math. Soc. \textbf{34} (2002), no.~2, 155--178.

\bibitem{Huybrechts2006}
D.~Huybrechts,
\emph{Fourier--Mukai Transforms in Algebraic Geometry},
Oxford Mathematical Monographs, The Clarendon Press, Oxford University Press, Oxford, 2006.

\bibitem{Kontsevich1994}
M.~Kontsevich,
\emph{Homological algebra of mirror symmetry},
Proceedings of the International Congress of Mathematicians, Vol.~1, 2 (Z\"urich, 1994), Birkh\"auser, Basel, 1995, pp.~120--139.

\bibitem{Seidel2008}
P.~Seidel,
\emph{Fukaya Categories and Picard--Lefschetz Theory},
Zurich Lectures in Advanced Mathematics, European Mathematical Society (EMS), Z\"urich, 2008.

\bibitem{Thomas2001}
R.~P.~Thomas,
\emph{Moment maps, Fourier--Mukai transforms and mirror symmetry in Tyurin rotations},
Adv. Theor. Math. Phys. \textbf{5} (2001), no.~5, 801--849.

\end{thebibliography}

\end{document}
